\section{Materials and methods}
\label{sec:method}

We assess where decarbonised hydrogen can compete with heat-pump electrification for
European building heat using a spatially resolved, bottom-up techno-economic model coupled
to a merit-order representation of the power sector. It is a chained simulation rather than
a system-wide optimisation, reconstructing useful heat demand region by region, splitting
it across eight technologies, pricing each on a levelised basis and then putting hydrogen
against its rivals in operating-cost dispatch tests, so hydrogen's share is set
exogenously and then tested, never co-optimised (Supplementary Section~S1.1). Coverage spans the EU-27, the United Kingdom and
Switzerland, and demand resolves to the third level of the Nomenclature of Territorial
Units for Statistics (NUTS3), 1{,}369 regions, aggregating to the country before the
economic tests apply. The model runs on the milestone years 2025, 2030, 2040 and 2050,
with 2025 the common anchor for demand, technology-share and cost trajectories alike
(Fig.~\ref{fig:arch}).

\begin{figure}[tbp]
 \centering
 \includegraphics[width=\textwidth]{P00_model_flow}
 \caption{\textbf{Model architecture.} Inputs, the demand reconstruction and the scenario
 draws feed three tests applied in order.}
 \label{fig:arch}
 \par\vspace{3pt}{\footnotesize\raggedright
 \textit{Note:} a technology that fails an earlier test is not rescued by a later one.
 The first two tests are cost comparisons at different time scales, annual and peak-hour,
 whereas the third asks whether the plant that wins the peak can be financed at all. The
 text below and the Supplementary Information give the data sources, the price paths, the
 adder bounds and the parameter ranges the boxes summarise.\par}
\end{figure}

The demand layer is built before the technology layer. A bottom-up strand reconstructs
residential heat from open building footprints and a harmonised archetype typology,
while an independent top-down baseline from the Hotmaps regional database
\citep{Hotmaps2020}, with census-derived building-type shares
\citep{Eurostat2021census,ONS2021TS044}, benchmarks it. The two agree to $-0.8$~per~cent across the aggregate and to within $\pm 25$~per~cent in 28
of the 29 countries. Section~\ref{sec:results:mix} gives the stricter vintage-matched
comparison (Supplementary Section~S1.2). National demand is then projected forward from 2025 on a multi-driver ratio identity,
with population growth and falling household size adding demand and envelope renovation
removing it. Boundary and terrain typologies come from the Eurostat GISCO service
\citep{EurostatGISCO2021}, and Fig.~\ref{fig:nuts3map} shows the demand field, whose
steep within-border gradient a country-resolution model averages away.

\begin{figure}[tbp]
 \centering
 \includegraphics[width=\textwidth]{P01_nuts3_map}
 \caption{\textbf{Useful residential heat demand by NUTS3 region, 2025.} From the bottom-up build.}
 \label{fig:nuts3map}
 \par\vspace{3pt}{\footnotesize\raggedright
 \textit{Note:} the colour scale is logarithmic over 0.3 to 12~TWh/yr, close to the 2nd
 and 98th percentiles of the mapped values, and the colourbar arrowheads mark the regions
 outside that range.\par}
\end{figure}

Each of the eight technologies, comprising gas, oil and biomass boilers, a resistance
heater, air-source and ground-source heat pumps, district heat and a hydrogen boiler, is
priced on a levelised cost of heat that annualises capital at the country
weighted-average cost of capital and adds fuel, fixed and variable operating costs, in
real EUR$_{2024}$ throughout (Supplementary Section~S1.3). Equipment capital is sized to peak heat demand and annualised over the heat that capacity delivers, on the reasoning that a household buys a boiler for its coldest day rather than its average one, so efficiency changes the fuel bill and not the purchase price. Annualising over rated capacity instead would divide that charge a second time by each technology's own efficiency, which is why this is the accounting choice that moves the base-load count furthest. Techno-economic parameters are anchored to the Danish Energy Agency Technology Catalogue
\citep{DEA2023technology} and cross-checked against JRC, IEA and IRENA cost data
\citep{JRC2023technology,IEA2022heatpumps,IRENA2022heatpumpcosts}. Fossil fuels carry the second emissions trading system for buildings and road transport
(ETS2) from its 2027 launch \citep{EuropeanCommission2024ETS2}, green hydrogen follows
the European Hydrogen Observatory central delivered-cost trajectory
\citep{EHO2024hydrogencosts}, and retail gas and electricity prices are Eurostat
first-half-2025 figures \citep{Eurostat2025electricityprices,Eurostat2025gasprices}.

A scope-1 inventory runs alongside the cost layer on the same technology split, charging
gas and oil fixed combustion factors \citep{Calvin2023Ipcc}, biomass and hydrogen zero
on the European renewable-fuel convention, and electricity and district heat
country-by-year intensities that each scenario's grid lever scales
\citep{EMBER2024electricity,EEA2024gridintensity}. Those paths are calibrated
trajectories rather than published projections, so the emissions levels inherit that
calibration whereas the scenario ranking does not (Supplementary Section~S1.4).
Furthermore, two last-mile network charges apply by default, one to each carrier, so that
neither is priced as though its delivery were free. The hydrogen boiler carries a coverage-weighted blend of conversion and new build at a median \euro{}13.4/MWh, and the heat pump a low- and medium-voltage reinforcement adder at a median \euro{}4.9/MWh, leaving the hydrogen route roughly 2.7 times the delivery charge of the electric one before any fuel cost enters.

Levelised cost compares technologies as base-load suppliers, an arena the review literature
treats as settled in the heat pump's favour
\citep{Rosenow2022PipeDream,Rosenow2024MetaReview,Billerbeck2023Race}. Hydrogen's narrower
case is as a peaker serving the cold-snap slice, modelled as a merit-order dispatch among
quick-start options whose SRMCs per MWh of useful heat are

\begin{equation}
\label{eq:srmc}
 c^{\mathrm{gas}}_{c}
 = \frac{p^{\mathrm{gas}}_{c} + \varphi_{y}\, e_{\mathrm{g}}}{\eta_{\mathrm{g}}},
 \qquad
 c^{\mathrm{H_2}}_{c}
 = \frac{p^{\mathrm{H_2}}_{c}}{\eta_{\mathrm{h}}} + \sigma_{c},
 \qquad
 c^{\mathrm{HP}}_{c}
 = \frac{\pi^{\mathrm{peak}}_{c}}{\mathrm{COP}^{\mathrm{cold}}_{c}},
\end{equation}

\noindent where $\varphi_{y}$ is the carbon price, $\eta$ the conversion efficiencies
and $\sigma_{c}$ the seasonal-storage adder, about \euro{}50/MWh-heat with domestic salt caverns and \euro{}100 without, the two ends spanning the cavern and above-ground rates reported in the bulk-storage literature \citep{Lord2014GeologicStorage,Papadias2021BulkStorage} and applied at the higher rate to the markets with no salt resource \citep{Caglayan2020caverns}. Cold-snap heat-pump COP is derated from the seasonal figure, $\mathrm{COP}^{\mathrm{cold}}_{c}=\max(0.60\,\mathrm{SCOP}_{c},1.8)$
following the low-temperature field measurements \citep{Gibb2023Coming}, which gives 1.80 to 2.84 across the 29
markets at a median of 2.04 in 2050. The scarcity electricity price
$\pi^{\mathrm{peak}}_{c}$ that prices the heat-pump branch is taken from the power
sector's own merit order, set by the marginal winter unit. A country counts as a hydrogen win where its hydrogen-ready turbine undercuts the gas turbine on SRMC, both machines held at the same efficiency.
Winning the dispatch is necessary but not sufficient, because an investor must recover the
peaker's capital from the few hundred hours a year it runs
(Eq.~\eqref{eq:recovery}). The question therefore shifts from an energy share to one of
firm capacity, screened consistently across the building winter peak, district heat and
the power sector. Symbols are listed in the nomenclature above.

Four multi-lever scenarios bracket the plausible range of European buildings
decarbonisation to 2050. \textbf{Current~Policies} freezes today's implemented measures;
\textbf{Stated~Policies} delivers announced policy, broadly consistent with the European
Commission long-term strategy \citep{EuropeanCommission2018LTS}; \textbf{Net~Zero}
matches the IEA Net~Zero scenario \citep{IEA2023NZE}; and \textbf{H2 Push} is a
deliberately hydrogen-favourable world along the lines of REPowerEU
\citep{REpowerEU2022}. Table~\ref{tab:levers_body} summarises the six exogenous levers each scenario sets, which separate them most sharply on renovation and on fossil retreat. Current~Policies renovates at 0.70 times the announced-policy rate and still leaves 40~per~cent of 2050 heat on fossil fuels, whereas Net~Zero renovates at 2.73 times that rate and leaves 5~per~cent. H2~Push is the hydrogen-favourable world rather than the most ambitious one, pairing a high carbon price and the rapid hydrogen-price path with a 1.64 times renovation rate and a 15~per~cent fossil share. Hydrogen adoption is set at 0, 5, 3 and 8~per~cent, so its ordering does not follow climate ambition. That setting is exogenous, and the dispatch screen tests it rather than sets it. Set against the dispatch-supported potential in the final row, which runs 0, 0.6, 7.5 and 9.0~per~cent, only Stated~Policies places adoption above what the screen supports, a gap Section~\ref{sec:results:merit} reports as a finding.

\begin{table}[htbp]
 \centering
 \caption{\textbf{Scenario lever settings, 2050.} Six exogenous levers per scenario, with
 the independently derived dispatch-supported potential beneath them.}
 \label{tab:levers_body}
 \begin{tabular}{lcccc}
 \toprule
 Lever & Current & Stated & Net~Zero & H2~Push \\
 & Policies & Policies & & \\
 \midrule
 Envelope rate (vs announced)\textsuperscript{a} & $\times$0.70 & $\times$1.00 & $\times$2.73 & $\times$1.64 \\
 Carbon-price path & Low & Central & High & High \\
 Hydrogen-price path & Stranded & Central & Central & Rapid \\
 Grid-decarbonisation speed\textsuperscript{a} & $\times$1.15 & $\times$1.00 & $\times$0.70 & $\times$1.00 \\
 Fossil share 2050 (ambition)\textsuperscript{c} & 40\% & 27\% & 5\% & 15\% \\
 H$_2$ adoption setting 2050\textsuperscript{b,c} & 0\% & 5\% & 3\% & 8\% \\
 \midrule
 \multicolumn{5}{l}{\textit{Derived independently, not a lever}} \\
 Dispatch-supported potential\textsuperscript{d} & 0\% & 0.6\% & 7.5\% & 9.0\% \\
 \bottomrule
 \end{tabular}
 \\[3pt]{\footnotesize\raggedright
 \textsuperscript{a} Multipliers on the announced-policy (Stated~Policies) central case.\\
 \textsuperscript{b} An exogenous adoption pathway, not an output of the dispatch screen.\\
 \textsuperscript{c} Share settings are Monte~Carlo means rather than hard caps, and the
 realised medians sit a little above them for the reason given in the text. Supplementary
 Table~S4 repeats the levers.\\
 \textsuperscript{d} The demand-weighted peak slice the symmetric screen supports, at the
 central scarcity premium and derate. Stated~Policies sets adoption above it, which
 Section~\ref{sec:results:merit} reports as a finding. The asymmetric accounting
 gives 0, 5.1, 7.9 and 9.3~per~cent (Supplementary Sections~S1.3 and~S1.6).\\
 \textit{Source: Authors' contribution.}\par}
\end{table}

\paragraph{Weather series for the power layer}
The weather series is the strongest assumption in the paper and every power-sector
number descends from it, so we state it here rather than leaving it to the Supplementary
Information. The power layer runs on synthetic hourly series, not reanalysis or metered
data, and on one realisation of a single year. Hourly heat demand is a seasonal profile
times a fixed daily shape, amplified over one six-day cold snap in mid-January; onshore
wind is an AR(1) process forced down to a lull over those same six days in all 29
countries at once; solar and the non-heat baseline load are deterministic. That
construction makes the event perfectly correlated across the study area, and the sizing
rule, which takes the fourth-highest residual hour of the year and so admits three hours
of lost load, is an order statistic on that one year rather than a loss-of-load
expectation over a distribution of weather years.

The firm-fleet requirement carries a band of about $\pm$9~per~cent from the assumed
event length alone, and a reader should treat it, the full-load hours and the cumulative
shortfall as conditional on this constructed year, and not as an adequacy assessment.

Three extensions strengthen the power-sector coupling. An explicit unit commitment
cross-checks the reduced-form peak-capacity cap, respecting minimum stable output, ramp
limits, minimum up and down times and an operating-reserve margin; a myopic
rolling-horizon test relaxes perfect foresight in the least-cost linear program; and an
endogenous Switzerland treatment couples incremental Swiss heat-pump load back to the EU
winter electricity price and credits the Swiss reservoir fleet as firm winter capacity.
Parametric uncertainty is characterised with 200 Monte Carlo draws per scenario,
sampling the envelope-renovation rate, fuel and hydrogen prices, equipment costs and
heat-pump performance under a cross-country shock correlation of 0.5 on renovation and 0.75 on fuel prices, set higher on prices because European gas and electricity markets clear against shared hubs whereas renovation policy is delivered nationally, and reported at the 10th, 50th and 90th percentiles. A Sobol screen ranks the four
demand drivers, and rank-regression coefficients the cost-side ones. Each extension's formulation, the sampling distributions and the
validation gates are in the Supplementary Information.

The model, the scenario configurations and the analysis code are released in a public
repository under an MIT licence, with retrieval scripts for the source datasets, so
every test reported below reproduces from the committed intermediate data; rebuilding
the regional demand layer from raw building footprints additionally requires the object
store named in the data-availability statement.
